\chapter{Manipulaci\'on de Datos (DML), Filtrado y Consultas Fundamentales}
\label{cap:dml_consultas_basicas}

\section{Marco Conceptual y Te\'orico (Curr\'iculo Universitario)}

El Lenguaje de Manipulaci\'on de Datos (DML) permite interactuar con el estado interno de las relaciones relacionales mediante operaciones de inserci\'on, actualizaci\'on, eliminaci\'on y consulta (DQL).

\subsection{L\'ogica de Tres Valores (3VL) y Tratamiento de NULLs}
En SQL, una condici\'on l\'ogica puede evaluarse en: \textbf{TRUE}, \textbf{FALSE} o \textbf{UNKNOWN}. Los valores \texttt{NULL} representan la ausencia de valor o valor desconocido:
\begin{itemize}
    \item La expresi\'on \texttt{columna = NULL} siempre eval\'ua a \textbf{UNKNOWN}.
    \item Se debe usar obligatoriamente el predicado formal \texttt{IS NULL} o \texttt{IS NOT NULL}.
    \item En ordenamientos, PostgreSQL coloca los \texttt{NULL} al final por defecto en orden descendente (\texttt{NULLS LAST}).
\end{itemize}

---

\section{Casos de Estudio y Pr\'actica Aplicada}

\begin{lstlisting}[style=sqlstyle, caption={Filtrado riguroso de pedidos con funciones escalares y fechas}]
-- Consulta anal\'itica de pedidos internacionales con c\'alculo de d\'ias de demora
SELECT
    order_id,
    customer_id,
    order_date,
    shipped_date,
    (shipped_date - order_date) AS days_to_ship,
    freight,
    COALESCE(ship_region, 'Sin Regi\'on Asignada') AS region_status
FROM orders
WHERE order_date >= '1997-01-01'::DATE
  AND shipped_date IS NOT NULL
  AND freight > 50.00
ORDER BY
    freight DESC,
    days_to_ship ASC
LIMIT 20;
\end{lstlisting}

---

\section{Taller de Consolidaci\'on del Cap\'itulo}

\begin{businessproblem}
El \'area de log\'istica de Northwind requiere identificar todos los productos descontinuados o con stock cr\'itico (menos de 10 unidades en almac\'en pero con pedidos pendientes de clientes), calculando el valor total monetario inmovilizado en inventario.
\end{businessproblem}

\subsection{Soluci\'on T\'ecnica en PostgreSQL}

\begin{lstlisting}[style=sqlstyle, caption={Consulta de evaluaci\'on de riesgo de stock e inventario}]
SELECT
    product_id,
    product_name,
    unit_price,
    units_in_stock,
    units_on_order,
    (unit_price * units_in_stock)::NUMERIC(12, 2) AS total_inventory_value,
    CASE
        WHEN discontinued = 1 THEN 'DESCONTINUADO'
        WHEN units_in_stock = 0 THEN 'ROTURA TOTAL DE STOCK'
        WHEN units_in_stock < 10 THEN 'ALERTA: STOCK CR\'iTICO'
        ELSE 'STOCK NORMAL'
    END AS inventory_health_status
FROM products
WHERE discontinued = 1
   OR (units_in_stock < 10 AND units_on_order > 0)
ORDER BY
    total_inventory_value DESC;
\end{lstlisting}

\begin{engtip}
\begin{itemize}
    \item \textbf{Casting Expl\'icito:} El uso de \texttt{::NUMERIC(12, 2)} previene errores de precisi\'on flotante (\textit{floating-point inaccuracies}) inherentes a tipos como \texttt{REAL} o \texttt{FLOAT}.
    \item \textbf{Funci\'on \texttt{COALESCE}:} Permite sustituir de forma segura valores desconocidos o nulos sin alterar el tipo de retorno.
\end{itemize}
\end{engtip}
